\documentclass[11pt]{article}
\usepackage{mathblatt}
\begin{document}
\weit
\typeout{PROBE-BEGIN trigonometrische-funktionen-e1-k1-s0-v3}
\begin{aufgabe}{\detokenize{trigonometrische-funktionen-e1-k1-s0-v3}}
\begin{teile}
\teil Probe

\einheitskreis[ohne]{65}
\end{teile}
\end{aufgabe}
\typeout{PROBE-END trigonometrische-funktionen-e1-k1-s0-v3}
\typeout{PROBE-BEGIN trigonometrische-funktionen-e1-k1-s2-v1}
\begin{aufgabe}{\detokenize{trigonometrische-funktionen-e1-k1-s2-v1}}
\begin{teile}
\teil Probe

\einheitskreis[ohne]{100}
\end{teile}
\end{aufgabe}
\typeout{PROBE-END trigonometrische-funktionen-e1-k1-s2-v1}
\typeout{PROBE-BEGIN trigonometrische-funktionen-e1-k1-s2-v2}
\begin{aufgabe}{\detokenize{trigonometrische-funktionen-e1-k1-s2-v2}}
\begin{teile}
\teil Probe

\einheitskreis[ohne]{235}
\end{teile}
\end{aufgabe}
\typeout{PROBE-END trigonometrische-funktionen-e1-k1-s2-v2}
\typeout{PROBE-BEGIN trigonometrische-funktionen-e1-k1-s2-v3}
\begin{aufgabe}{\detokenize{trigonometrische-funktionen-e1-k1-s2-v3}}
\begin{teile}
\teil Probe

\einheitskreis[ohne]{290}
\end{teile}
\end{aufgabe}
\typeout{PROBE-END trigonometrische-funktionen-e1-k1-s2-v3}
\typeout{PROBE-BEGIN trigonometrische-funktionen-e1-k1-s5-v1}
\begin{aufgabe}{\detokenize{trigonometrische-funktionen-e1-k1-s5-v1}}
\begin{teile}
\teil Probe

\einheitskreis{50}
\end{teile}
\end{aufgabe}
\typeout{PROBE-END trigonometrische-funktionen-e1-k1-s5-v1}
\end{document}
